\documentclass[b5paper, 10pt]{ltjsarticle}
\usepackage{custom4-default}
\usepackage{gb4e}\noautomath
\usetikzlibrary{arrows.meta}
\newcommand{\ko}[1]{\textsf{#1}}
\newcommand{\sentence}[1]{$\langle$#1$\rangle$}
\newcommand{\kox}[1]{\textsf{[#1]}}
\usepackage[safe]{tipa}

\begin{document}
% 全11節の試稿．執筆規約・参照先は editorial-plan.md を参照．
\title{コノメノ素描}
\author{}
\date{}
\maketitle
\setcounter{tocdepth}{1}
\tableofcontents

\section{コノメノとは}
コノメノ（\ko{konomeno}）は，琴粥が制作している人工言語である．名前は\ko{kono}「概念」と\ko{meno}「法則」からなり，「概念の法則を記述する言語」という関心を表している．語をどう並べるかに加え，何を一つのものとして扱い，どのような関係で結び，そこから何を言えるかを考えながら作っている．

例えば「大きい」と「大きさがある」は同じではない．小さいものにも大きさはある．また，「カゴの中にりんごがある」と「りんごはカゴの一部である」も違う．日常的には気に留めず使っている表現でも，その違いを明示しようとすると，量や部分や所在をどう捉えるかという問題になる．コノメノでは，こうした概念のモデルを語彙と文法に組み込んでいく．

\subsection{何を作りたいか}
音や文字には，芸術言語としての関心が強い．音の色や手触り，動きの印象をもとに語を作り，語彙全体に固有の表情を持たせたい．文法と意味には，工学言語としての関心が強い．文の意味を形式的に書き，語の意味に密着した推論を扱えるようにしたい．ただし，理論自体も創作作品として作っている．少ない原理で多くを説明することや，さまざまなものを分類して立ち位置を明確にすることにも，美意識がある．

モデル化の仕方は一通りではない．何を細かく区別し，何を一つにまとめるかによって，得意な表現も変わる．そのため，意味論やオントロジーの複数のモデルを，一つの言語の中で使えるようにすることも目標にしている．これらは制作の方向であり，あらゆる日常表現をすでに形式化できるという意味ではない．

\subsection{口語と文語}
人が話したり書いたりするための言語が\emph{口語コノメノ}である．本書では，口語の短い文から読み方を学ぶ．口語には，文脈に頼って対象を指したり，関係の向きや項を省いたりする工夫がある．短く言える分だけ，聞き手が解釈を補うこともある．

一方，\emph{文語コノメノ}は，文法や意味を明示的に扱う形式言語である．基本的には記号で書き，発話のための音韻を持たない．口語の省略を解消した形を，本書でそのまま別の文体として文語と呼ぶわけではない．例文の後に図や式を添えるときには，何をどの関係で結んでいたかを見えるようにするために使う．

\subsection{本書の読み方}
論理学や言語学の前提知識は仮定しない．まず例文を読み，次にその形を作る規則を説明する．必要な用語はそのつど導入するので，最初から分類名をすべて覚える必要はない．同じりんごや食事の例を使いながら，文を少しずつ長くしていく．

本書は，短い文を読み書きする入口であるとともに，詳しい仕様書『コノメノ詳説』を読むための見取り図でもある．LLMと共同執筆するときにも，装置の用途や制作意図を共有する足場として使う．現行の定義と未決事項は『詳説』本文に置き，本書では理解に必要な範囲を選んで説明する．定義を確認したいときの案内は最後の節にまとめた．

\section{音と綴り}
本書ではラテン文字による綴りを使う．コノメノにはマリータスユという独自の表音文字もあるが，まずは以下の綴りで例文を読めるようにしよう．ここでの発音説明は，文法の例を読むための入口である．

\subsection{母音と子音}
母音字は\ko{a, e, i, o, u, v, y}の7個である．\ko{a, e, i, o}はおおむね日本語のア・エ・イ・オを手掛かりにできる．残る3個は，次のように区別する．角括弧内は国際音声記号による発音である．
\begin{center}
\begin{tabular}{clp{0.60\linewidth}}
    \toprule
    綴り & 発音 & 口の形の目安 \\
    \midrule
    \ko{u} & \textipa{[9]}$\sim$\textipa{[W]} & 唇を丸めない．日本語のウよりも唇をすぼめず，曖昧な母音に近い範囲も持つ．\\
    \ko{v} & \textipa{[u]} & 唇を丸めて発音するウ．子音字ではない．\\
    \ko{y} & \textipa{[y]} & イの舌の位置を保ちながら唇を丸める．\\
    \bottomrule
\end{tabular}
\end{center}
子音には\ko{p, b, t, d, k, g, f, s, z, h, m, n, l, w, j}のほか，\ko{c, zc, ts, tc, kh, ng}がある．\ko{j}はヤ行の子音に当たり，\ko{ja}は「ヤ」に近い．\ko{l}は日本語のラ行に近い音から舌先を付ける[l]までを含み，rとlの対立はない．\ko{b}には摩擦音[v]としての実現もある．

綴りから読み違えやすいものを挙げる．
\begin{center}
\begin{tabular}{clp{0.60\linewidth}}
    \toprule
    綴り & 発音の例 & 目安 \\
    \midrule
    \ko{c} & \textipa{[C]}$\sim$\textipa{[S]} & シの子音に近い．\\
    \ko{zc} & \textipa{[\t{d\textctz}]} & ジの子音に近い．\\
    \ko{ts} & \textipa{[ts]} & ツの子音．後ろに母音ウを足さない．\\
    \ko{tc} & \textipa{[\t{tC}]}$\sim$\textipa{[\t{tS}]} & チの子音に近い．\\
    \ko{kh} & \textipa{[\c{c}]}$\sim$\textipa{[x]}$\sim$\textipa{[X]} & ヒの子音から，舌の奥で摩擦させる音まで．\\
    \bottomrule
\end{tabular}
\end{center}
二文字の\ko{tc}や\ko{kh}は一つの子音を表す．例えば人名\ko{komatci}は\ko{ko.ma.tci}と区切る．\ko{wa.ka}や\ko{ka.lam}の点も，ここでは説明用の音節境界であり，通常の綴りには入れない．子音で終わる音節に，日本語の仮名に合わせて母音を付け足さないようにする．

\subsection{綴りと実際の発音}
母音が続いていても，綴りを一字ずつ読むとは限らない．例えば\ko{ae}・\ko{ea}は\textipa{[\ae:]}，\ko{oy}は\textipa{[\o:]}，\ko{av}・\ko{ov}は\textipa{[o:]}となる．同じ母音を重ねた綴りは長い母音を表す．また，語末の単独の破裂音は，喉で息を止める\textipa{[P]}として発音する．

語と語の間には空白を置き，発音でも語境界を越えて音節を組み直さない．接辞を付けるときは語の内部の結合規則が働くので，単に文字列を連結するとは限らない．文法の説明に出る\ko{+ik}のようなプラス記号は，接辞を付ける操作の記号で，発音しない．実際の語形はその例ごとに示す．

\subsection{高さと語の印象}
コノメノは高低アクセントを持つ．現行の暫定規則では，音節の中で下降する型，音節と音節の間で下降する型，下降しない型を区別し，型を語ごとに指定する．綴りだけからすべての語のアクセントを一意に決めることはしない．

音節内で下降する型には，個別指定がなければ，単音節語ではその音節に，多音節語では語末を除いて左から最初の重音節に下降を置くという規則がある．長母音・二重母音を含む音節，子音で終わる音節，二つの子音で始まる音節は重い．多音節語で該当する音節がなければ後ろから2番目を選ぶ．下降を実現するために，母音を長くする場合もある．他の型や個々の語の指定まで，この規則で上書きしない．

造語では，発音できることに加えて，その語をどんな印象にしたいかを考える．作者の印象では，例えば\ko{y}には赤く柔らかな感じ，\ko{l}には細く揺れる感じがある．より長い音列にも，\ko{la}に食べ物，\ko{fon}に袋といった連想が育っている．こうした音列を素材にするが，音を聞けば語義が自動的に分かるという符号表ではない．連想は語形を選ぶ手掛かりで，出来上がった語の定義は別に定める．

独自文字マリータスユでは，子音と母音の字形を音節ごとに組み合わせる．読み書きの効率とともに，抽象絵画のような文字の美しさや可愛らしさを重視している．字形とその組み立て方は『詳説』の文字章で扱う．

\section{単関係文と指示}
コノメノでは，物事の間の関係を一つずつ書き出して文を作る．まずは，二つの対象を一つの関係で結ぶ\emph{単関係文}を見てみよう．例文の\sentence{山括弧}には日本語訳を書く．

\subsection{二つの対象を結ぶ}
小町が自己紹介する場面を考える．\ko{waka}は「私」，\ko{komatci}は人名の「小町」，\ko{kast}は「同じ対象である」という関係を表す．これらを次の順に並べる．
\begin{exe}
    \ex \ko{waka kast komatci}\hfill \sentence{私は小町である}
\end{exe}
\ko{waka}と\ko{komatci}はそれぞれ対象を指し，\ko{kast}はその二つの指し方が同じ人を指すと述べている．このように対象を指す表現を\emph{項}，二つの項を結ぶ関係を\emph{二項関係}と呼ぶ．単関係文の基本形は「項・関係・項」である．関係に渡す二つの項を，順に第1項・第2項と呼ぶこともある．

「私」が誰であるかは話者によって変わる．小町が言えば小町を，別の人が言えばその人を指す．これに対して，固有名\ko{komatci}は話者が変わっても小町を指す．\ko{non}「あなた」も\ko{waka}と同じく発話の場面から指す相手が決まる語で，このような語を\emph{指標詞}という．

単関係文は，関係語を替えることでいろいろな内容を表せる．例えば\ko{micto}は「第1項を第2項が所有する」という関係である．ただし，何を所有しているかを言うためには，まず普通の名詞で対象を指す方法が要る．

\subsection{概念から対象を指す}
机の上のりんごを指して，「これは私のものだ」と言いたいとする．「りんご」は\ko{kalam}なので，次のように言える．
\begin{exe}
    \ex \ko{kalam micto waka}\hfill \sentence{そのりんごは私のものだ}
\end{exe}
\ko{kalam}は，語としては「りんご」という\emph{概念}を表す．しかし，この文で所有されているのは「りんご」という概念そのものでも，世の中のりんご全部でもない．いま話題にしている一つのりんごである．概念語を項の位置に置くと，このように発話の場面からその概念に当てはまる対象を取る．

この働きは\emph{冠詞}\ko{fe}で明示できる．冠詞とは，概念からどのように項を作るかなどを指定する語である．ここでは\ko{fe}だけを使う．
\begin{exe}
    \ex \ko{fe kalam micto waka}\hfill \sentence{そのりんごは私のものだ}
\end{exe}
\ko{fe kalam}は「場面から取ったりんご」を指す．冠詞のない\ko{kalam}も，この項位置では同じ働きをする．以後の例文でも，この省略を使う．固有名や\ko{waka}は初めから対象を指すので，概念語と同じ仕方で\ko{fe}を補う必要はない．

訳の「その」は，すでに文章に登場したことを要求しない．目の前にあるりんごについて初めて話すときにも使える．反対に，りんごがいくつもあってどれを指すか分からなければ，\ko{fe}を発音するだけで解決するわけではない．指差しや説明を足して，どの対象について述べているかを伝える必要がある．

\begin{remark}
    概念と，その概念に当てはまる対象を分けるのは，「どのりんごか」「すべてのりんごか」といった違いを文の意味に反映させるためである．冠詞を替えると何が変わるかは「量化・否定・複数」の節で扱う．
\end{remark}

\subsection{関係の向きと使い分け}
\ko{micto}の第1項は所有されるもの，第2項は所有者である．したがって，\ko{kalam micto waka}は，日本語の語順に寄せれば「私はりんごを所有する」と訳せるが，コノメノの語順は「りんご・所有される・私」である．関係語を覚えるときは，訳語と一緒にこの向きを覚える．

次の二つも同じ「項・関係・項」の形をしている．\ko{misuy}は「指」，\ko{kina}は「手」，\ko{folka}は「カゴ」である．
\begin{exe}
    \ex \ko{misuy hanov kina}\hfill \sentence{その指はその手の一部だ}
    \ex \ko{kalam litam folka}\hfill \sentence{そのりんごはそのカゴの中にある}
\end{exe}
\ko{hanov}は「第1項は第2項の物理的な部分である」，\ko{litam}は「第1項は第2項の中にある」を表す．どちらも二つのものを結ぶが，カゴの中のりんごは，それだけではカゴの一部ではない．また，カゴに入っていることから所有者が誰かも決まらない．

日本語では「私のりんご」「手の指」「カゴのりんご」のように，いずれも「の」で結べる．コノメノでは，どの関係について述べているかを語で選び分ける．こうした区別を作るのは，文から何が分かり，何までは分からないかを整理するためである．単語を増やす作業と，関係の性質を調べる作業が結び付いている．

\begin{remark}
    \ko{hanov}の例では，指や手を物理的な部分の集まりとして見ている．『詳説』では，手のような一つのまとまりを持つものと，物理的な部分の集まりを区別し，この見方の切替えを省略として扱う．ここでは日常的な形を使うが，両者が同じ種類の対象というわけではない．
\end{remark}

\subsection{「人である」と「その人である」}
「小町は人である」は，人という概念に小町が当てはまることを述べる．この場合は，\ko{coto}「人」を小町の後に直接置く．
\begin{exe}
    \ex \ko{komatci coto}\hfill \sentence{小町は人である}
    \ex \ko{komatci kast coto}\hfill \sentence{小町はその人である}
\end{exe}
最初の文の\ko{coto}は，「人である」という条件として働く．一つの項について成り立つ条件を\emph{一項述語}と呼び，「項・一項述語」の形を\emph{所属文}と呼ぶ．後の文では，\ko{kast}の右側が項の位置なので，\ko{coto}は場面から取った人を指す．小町とその人が同一だ，という意味になる．

例えば，小町が人間だと説明したいときと，写真の中の人が小町だと教えたいときでは，言いたいことが違う．日本語の「である」を一律に\ko{kast}に置き換えると，この違いが消えてしまう．所属文では概念を条件として使い，同一性の文では二つの項を\ko{kast}で結ぶ．

\subsection{すでに出た対象を指す}
一度話題に出した対象をもう一度指すことを\emph{照応}という．場面からりんごを取る\ko{fe kalam}と，すでに出たりんごを受ける表現は区別する．例えば，りんごを示してから，その所有者について述べるなら次のように言える．
\begin{exe}
    \ex \ko{kalam, seta micto waka}\hfill \sentence{りんごだ．それは私のものだ}
\end{exe}
最初の\ko{kalam}は，項を一つ置いて対象を話題に出す\emph{提示文}である．\ko{seta}「それ」は，最も新しく言及された対象を受ける．ここではりんごを受け，その後で\ko{waka}が私を指す．コンマはここでは，二つの内容を続けて述べる区切りとして使っている．

同じ概念語を繰り返すだけでは，再び同じ対象を指すことは保証されない．\ko{kalam, kalam}は，そのつど場面からりんごを取る表現になる．二つが同じりんごであるとも，別のりんごであるとも，この反復だけでは言っていない．同じ対象を受けたいときには，\ko{seta}のような指示詞を使う．固有名の\ko{komatci}は，繰り返しても小町を指す．

\section{動詞文と一筆書き}
「私はりんごを食べる」は，次のように言う．この文を，前節の単関係文から組み立ててみよう．
\begin{exe}
    \ex \ko{waka ja hamin ke kalam}\hfill \sentence{私はりんごを食べる}
\end{exe}
この例では，ある食事と，そこで食べる一つのりんごについて述べている．\ko{hamin}は「食べる」を表す語だが，この位置では\ko{kalam}と同じく項として使われ，その食べる出来事を指す．

\subsection{出来事と参与者}
「私」と「りんご」を直接結ぶ代わりに，食べる出来事を一つ取り出し，両方をその出来事に結ぶ．出来事に関わる人やものを\emph{参与者}，どのように関わるかを表す関係を\emph{意味役割}という．ここでは次の二つを使う．
\begin{itemize}
    \item \ko{ja}：第1項が第2項の動作主である．この例では，私が食べる行為者であることを表す．
    \item \ko{ke}：第1項が第2項の被破壊物である．この例では，りんごの存在段階が，食べる出来事の終わりで終わることを表す．
\end{itemize}
どちらも第2項に出来事を取る二項関係である．それぞれの向きで書けば，次の二つの単関係文になる．
\begin{exe}
    \ex \ko{waka ja hamin}\hfill \sentence{私は食べる}
    \ex \ko{kalam ke hamin}\hfill \sentence{そのりんごが食べられる}
\end{exe}
\ko{hamin}の項位置での読みも，省略された\ko{fe}による．「食べるという概念」ではなく，場面から取った食べる出来事を指している．ただし，この二つをそのまま続けるだけでは，同じ出来事を指すことが保証されない．一つの食事について述べるには，後の\ko{hamin}が前の出来事を受けるようにする必要がある．

\begin{remark}
    \ko{ke}は日本語の「を」一般に対応する語ではない．「りんごを見る」のりんごには，見られるだけで存在が終わるという意味はない．使う意味役割は，出来事をどう捉える語かに応じて選ぶ．食事中のりんごがいつまで同じりんごか，どの部分まで食べたかも，\ko{ke}だけでは決まらない．
\end{remark}
\begin{remark}
    動作主\ko{ja}の定義は，意図の形式化に未整備の部分があるため暫定である．ここでは『詳説』の食事の例に従って用いる．
\end{remark}

\subsection{同じ出来事を共有する}
\ko{seta}は最も新しく言及された対象を指した．名詞を添えて「その食事」と言いたい場合は\ko{sen hamin}を使える．二つ前の対象なら，それぞれ\ko{tata}，\ko{tan hamin}になる．

これを使って，二つの単関係文を同じ食事に結び付ける．
\begin{exe}
    \ex \ko{waka ja hamin, kalam ke tan hamin}
    \glt \sentence{私は食べる．その食事で，そのりんごが食べられる}
\end{exe}
\ko{tan hamin}を読む時点では，りんごが一つ前，食事が二つ前，私が三つ前である．だから\ko{sen}ではなく\ko{tan}を使う．数えるのは単語数ではなく，対象に言及した新しい順であり，関係語\ko{ja}や\ko{ke}は数えない．また，\ko{hamin}を添えても，食事だけに絞って数えるわけではない．

参照した対象は，そのつど最も新しく言及された対象になる．上の文末の直後なら，同じ食事は\ko{sen hamin}または\ko{seta}で受けられる．指示詞の番号は，最初に登場したときの番号として固定されていない．

この文の内容を図にすると，次のようになる．矢印は，関係の第1項から第2項へ向けて描く．
\begin{center}
\begin{tikzpicture}[>=Stealth]
    \node (i) at (0,0) {私};
    \node[align=center] (e) at (3.1,0) {食べる出来事};
    \node (a) at (6.2,0) {りんご};
    \draw[->] (i) -- node[above]{\ko{ja}} (e);
    \draw[->] (a) -- node[above]{\ko{ke}} (e);
\end{tikzpicture}
\end{center}
出来事を一つの対象として置くと，食べる人と食べられるものが，同じ出来事にどう関わるかを書ける．さらに別の参与者や時間について述べるときにも，その出来事に結び付けられる．文の骨格を関係の図として考えるのは，こうして必要な情報を足していくためである．

\subsection{関係を逆向きに読む}
この図を「私，食べる出来事，りんご」の順になぞると，\ko{ja}の矢印には沿って進むが，\ko{ke}の矢印には逆らって進む．関係の第1項と第2項を交換する操作を\emph{転置}といい，\ko{me}を付けて表す．
\begin{exe}
    \ex\begin{xlist}
        \ex \ko{kalam ke hamin}
        \ex \ko{hamin meke kalam}
    \end{xlist}
    \glt \sentence{その食事で，そのりんごが食べられる}
\end{exe}
上の二つは，同じりんごと食事について言えば同じ内容である．\ko{meke}は\ko{ke}を逆向きにした関係で，「第2項が第1項の被破壊物である」を表す．受動態のようにも見えるが，動詞だけでなく二項関係に使える操作である．

\ko{hamin}は出来事，\ko{kalam}はりんごなので，この例では\ko{ke}のどちら側に置く対象かが分かる．そのため，会話では転置を表す\ko{me}を省いて\ko{hamin ke kalam}とも言える．このように，関係が要求する対象の種類に照らして向きが決まる場合には，転置を補って読む．両方向が可能な関係まで，聞き手が話者の意図に合わせて自由に裏返すわけではない．

\subsection{一筆でつなぐ}
二つ目の単関係文を転置すると，同じ食事が一つ目の末尾と二つ目の先頭に並ぶ．今度は二つ目の文の先頭で食事を受けるので，\ko{sen hamin}が使える．
\begin{exe}
    \ex\begin{xlist}
        \ex \ko{waka ja hamin, sen hamin meke kalam}\hfill 対象を再指示
        \ex \ko{waka ja hamin meke kalam}\hfill 共有する項を重ねる
        \ex \ko{waka ja hamin ke kalam}\hfill 転置の表示も省略
    \end{xlist}
    \glt \sentence{私はりんごを食べる}
\end{exe}
同じ食事を指す項を重ねて一度だけ言うと，「項・関係・項・関係・項」の形になる．これを\emph{関係連鎖文}と呼ぶ．最後の形が，節の冒頭で示した文である．二つの関係は残っていて，\ko{hamin}がその両方につながっている．

図を一筆でなぞるように読めるので，この組み立て方を\emph{一筆書き}という．図の同じ部分を重ねて発音せず，対象から対象へ移りながら関係を述べられる．逆に，りんごから出発してもよい．
\begin{exe}
    \ex \ko{kalam ke hamin ja waka}\hfill \sentence{そのりんごを私が食べる}
\end{exe}
この場合は\ko{ke}が元の向き，\ko{ja}が逆向きである．語順を替えるには，図のどこからどの道を通るかを考える．単語を自由に並べ替えてよいわけではない．

例えば，食事の文の末尾へ\ko{litam folka}を足すとどうなるだろうか．
\begin{exe}
    \ex \ko{waka ja hamin ke kalam litam folka}
    \glt \sentence{私は，そのカゴの中にあるりんごを食べる}
\end{exe}
追加された\ko{litam}が結ぶのは，直前のりんごとカゴである．食べる出来事をカゴの中に置く文にはならない．また，「食べる出来事，りんご，私」の順に\ko{ja}を挟めば，りんごと私を動作主関係で結ぶことになってしまう．一筆でつながらない内容は，いくつかに分け，同じ対象への指示で結べばよい．

\subsection{省略と時制}
共有する項を連鎖の中で重ねるほか，先に言及した対象を指す項を省くこともできる．省略位置を，解説では\ko{\_}で示す．これは発音する語ではない．
\begin{exe}
    \ex \ko{waka ja hamin, kalam ke \_}
    \glt \sentence{私は食べる．そのりんごが，その食事で食べられる}
\end{exe}
空いた位置には\ko{ke}の第2項となる出来事が必要なので，先に出た食事を受ける．省略は，直前の単語を機械的に繰り返すことではない．関係に合う種類の対象を先行する言及から探す．同じ関係のもう一方の項と同じ対象を省略で受けることもできないので，自分自身との関係を述べるときは指示詞などで明示する．

短く話せるように省略を許し，その解釈には対象の種類や言及順を使う．これが口語の設計上の工夫である．ただし候補が一つに決まらない場合もある．そのときは，指示詞で対象を選ぶか，文を分け直して伝える．

最後に，出来事の時間を指定してみよう．過去は\ko{+ik}，未来は\ko{+as}を出来事の語に付ける．この例では\ko{hamin}が\ko{haminik}，\ko{haminas}になる．
\begin{exe}
    \ex \ko{waka ja haminik ke kalam}\hfill \sentence{私はそのりんごを食べた}
    \ex \ko{waka ja haminas ke kalam}\hfill \sentence{私はそのりんごを食べるだろう}
\end{exe}
現在の\ko{+en}もあるが，発話時が出来事の進行中に入るという現行の定義は暫定である．時制を付けない形を，それだけで「現在の習慣」を表す形とはしない．ここまでの\ko{hamin}の例では，いつの食事かは場面に委ねていた．

\section{関係節と修飾}
りんごを場面から指すだけでなく，「私のりんご」「カゴに入っているりんご」のように条件を添えて指したいことがある．コノメノでは，条件を述べる文から，その条件を満たすものを表す句を作る．これが\emph{関係節}である．

\subsection{条件を満たすもの}
\ko{kalam micto waka}は「そのりんごは私のものだ」だった．これを「私が所有するりんご」という条件にするには，\ko{kalam}に冠詞\ko{le}を付け，文を\ko{les}と\ko{fol}で囲む．
\begin{exe}
    \ex \ko{les le kalam micto waka fol}\hfill \sentence{私が所有するりんご}
\end{exe}
\ko{le kalam}は，場面から特定のりんごを選ぶ位置ではない．りんごをいろいろ当てはめて，後の条件を満たすか調べる位置である．\ko{les}は条件を作る範囲の始まり，\ko{fol}はその終わりを表す．この例なら，私が所有するりんごが条件に当てはまる．

条件に当てはまるものを一括して扱うとき，その集まりを\emph{クラス}と呼ぶ．ここで作ったのは「私が所有するりんご」のクラスであり，その中の一つのりんごでも，それらを一つにまとめた集合個体でもない．空欄に入れた対象について，条件が成り立つかどうかを答えるものだと考えるとよい．

関係節を作る位置は文頭に限らない．動詞文でも，条件を付けたい位置に\ko{le}を置く．
\begin{exe}
    \ex \ko{les le coto ja hamin ke kalam fol}
    \glt \sentence{そのりんごを食べる人}
    \ex \ko{les waka ja hamin ke le kalam fol}
    \glt \sentence{私が食べるりんご}
\end{exe}
最初は人の位置，後はりんごの位置を変えながら条件を調べる．\ko{le}を付けていない\ko{hamin}などの裸項は，これまでと同じく場面から取る．したがって，この形だけで「あらゆる食事で」「いつも食べる」という意味になるわけではない．

\subsection{条件から一つを指す}
「私が所有するりんご」という条件を作っただけでは，どのりんごについて次の文を述べるかは決まっていない．その条件に当てはまる対象を場面から取るには，出来上がった句の外側に\ko{fe}を付ける．
\begin{exe}
    \ex \ko{fe les le kalam micto waka fol litam folka}
    \glt \sentence{その私のりんごは，そのカゴの中にある}
\end{exe}
この文の最初の項は，\ko{fe}から\ko{fol}までである．\ko{les le kalam micto waka fol}が条件を作り，外側の\ko{fe}がその条件に当てはまるりんごを指す．その後，\ko{litam}がそのりんごとカゴを結ぶ．

\ko{le}と\ko{fe}は逆向きの仕事をしているように見えるが，何でも打ち消し合うわけではない．「私が所有する」という条件は残るし，その条件に当てはまるりんごがなければ，対象を指せない．条件を作ることと，条件を満たすものを指すことを分けておくと，「そのようなものがある」「そのようなものはすべて」とも組み合わせられる．それが次節の量化につながる．

\subsection{括弧が必要になるところ}
\ko{les}と\ko{fol}は，紙の上の括弧に当たる働きを，発音できる語で行っている．関係節に何を含めるかが変われば，作る条件も変わる．例えば，所有者についての条件と，所在についての条件は，次のように書き分ける．
\begin{exe}
    \ex \ko{les le kalam micto waka fol}
    \glt \sentence{私が所有するりんご}
    \ex \ko{les le kalam litam folka fol}
    \glt \sentence{そのカゴの中にあるりんご}
\end{exe}
どちらもりんごのクラスを作るが，最初は所有者，後は所在を条件にしている．二つの条件を同じりんごについて述べたいなら，内部でもその対象を共有する必要がある．単に文を二つ足して，別々のりんごを指してしまってはいけない．

本書では，新しい操作を導入する例の括弧を明示する．短い会話では省かれる箇所もあるが，まずは，項の範囲と条件全体の範囲を自分で区切れるようにしておくとよい．関係節を作るために新しい語彙を覚え直す必要はなく，すでに作れる文のどの位置を取り出すかを選べばよい．

\section{量化・否定・複数}
「そのりんごは私のものだ」「私のりんごがある」「りんごはすべて私のものだ」では，同じ関係を使っていても主張する範囲が違う．この違いを，項に付ける冠詞と，その働く範囲によって表す．

\subsection{あるものとすべてのもの}
\ko{kon}は「あるものが」，\ko{san}は「すべてのものが」に当たる冠詞で，基本の形では概念語の後に付ける．\ko{fe}や\ko{le}と違って後置が基本である．
\begin{exe}
    \ex \ko{kalam micto waka}\hfill \sentence{そのりんごは私のものだ}
    \ex \ko{kalam kon micto waka}\hfill \sentence{私が所有するりんごがある}
    \ex \ko{kalam san micto waka}\hfill \sentence{りんごはすべて私のものだ}
\end{exe}
\ko{kon}の文は，条件を満たすりんごが少なくとも一つあると言う．どれか一つを指差してからでなくても述べられるし，一つだけだとも言っていない．\ko{san}の文は，考えている範囲のりんごの一つ一つについて，同じ条件を要求する．\ko{kon}による\emph{存在量化}と，\ko{san}による\emph{全称量化}を合わせて\emph{量化}という．

「すべて」の範囲を狭めたければ，前節の関係節が使える．
\begin{exe}
    \ex \ko{les le kalam litam folka fol san micto waka}
    \glt \sentence{そのカゴの中にあるりんごは，すべて私のものだ}
\end{exe}
\ko{san}の前にある関係節が，調べる対象をカゴの中のりんごに絞る．外側の文は，その一つ一つを私が所有すると述べる．範囲を狭める条件と，その範囲についての主張を分けて作れる．

\begin{remark}
    全称量化だけでは，範囲に対象が存在することは要求しない．カゴの中にりんごがなければ，私のものでないりんごもないので，上の全称文は成り立つ．りんごが入っていることも述べたいなら，存在の主張を加える．
\end{remark}

\subsection{量化の順序}
「すべてのりんごには所有者がいる」と「すべてのりんごを所有する人がいる」は違う．前者ではりんごごとに所有者が違ってよいが，後者では一人の所有者が全りんごについて条件を満たす．この違いを，量化の\emph{作用域}という．作用域は，どこまでを一まとまりの条件として調べるかを表す．

二つの冠詞を区別して対応付けるため，ここでは\ko{sanoy}と\ko{konal}を使う．末尾の\ko{oy}と\ko{al}はそれぞれラベル1・2であり，「一つ」「二つ」を数えているわけではない．対応する括弧の端は\ko{sasoy}と\ko{kosal}になる．
\begin{exe}
    \ex \ko{fol kalam sanoy micto coto konal kosal}
    \glt \sentence{どのりんごにも，それを所有する人がいる}
    \ex \ko{fol kalam sanoy micto coto konal sasoy}
    \glt \sentence{すべてのりんごを所有する人がいる}
\end{exe}
最初の文では，\ko{fol}から\ko{kosal}までが存在量化の範囲である．外側の全称量化がりんごを一つずつ調べ，そのたびに所有する人がいるかを見る．後の文では，内側の\ko{sasoy}が全称量化を閉じる．外側で人を一人選び，その人がすべてのりんごを所有するかを見る．一番外側の括弧は省略している．

括弧を明示しない\ko{kalam san micto coto kon}は，現行の規約では左の冠詞ほど広い作用域を持ち，最初の読みになる．ただし，順序が問題になる文では括弧を使えば，読み手が省略規則を補う負担を減らせる．

\subsection{否定の範囲}
文の否定は\ko{meide}と\ko{tos}で囲んで表す．
\begin{exe}
    \ex \ko{meide kalam micto waka tos}\hfill \sentence{そのりんごは私のものではない}
\end{exe}
否定と量化の順序も区別する必要がある．以下では\ko{san}に対応する括弧端形\ko{sas}を使い，両方の括弧を明示する．一つの全称量化だけなので，ラベルは省いている．
\begin{exe}
    \ex \ko{meide fol kalam san micto waka sas tos}
    \glt \sentence{りんごがすべて私のもの，というわけではない}
    \ex \ko{fol meide kalam san micto waka tos sas}
    \glt \sentence{どのりんごも私のものではない}
\end{exe}
最初は「全部が私のもの」を否定しているので，私のものでないりんごがあればよい．私のりんごも混じっていてよい．後は一つ一つのりんごについて「私のものではない」を要求する．一個だけ私のりんごが混じっていても成り立たない．

一語の述語の否定を表す\ko{be+}もあるが，量化を含む文全体の否定とは仕事が違う．本書では，どこを否定したかを見失わないように，文否定の括弧を使う．

\subsection{クラスと集合個体}
関係節で作ったクラスは，対象が条件に当てはまるかを表していた．一方，「小町と私を要素に持つ集合」のように，集まり自体を一つの対象として扱うこともできる．それを\emph{集合個体}と呼ぶ．集合の要素である関係は\ko{litoi}で表す．
\begin{exe}
    \ex \ko{fonoi waka komatci tos}\hfill \sentence{私と小町からなる集合}
    \ex \ko{waka litoi fonoi waka komatci tos}
    \glt \sentence{私は，私と小町からなる集合の要素だ}
\end{exe}
\ko{fonoi}と\ko{tos}の間に要素を並べる．順序は問わず，同じ対象を重ねて挙げても要素数は増えない．小町自身が\ko{waka}と言っている場合なら，この集合は一人だけからなる．

概念の外延を集合個体にする接辞は\ko{+im}である．例えば\ko{cotoim}は人の集合を指す．小町という個体を一つだけ入れた集合は\ko{komatcykof}で，個体から単元集合を作る\ko{+ykof}を使う．小町自身と，小町だけからなる集合は同じではない．

\begin{remark}
    どのクラスも\ko{+im}で集合個体にできるわけではない．本書で扱う人やりんごの外延は集合にできる場合を考えるが，例えば「すべての集合」のクラスを一つの集合にすることは認めない．条件を作れることと，その外延を一つの対象にできることを分けている．
\end{remark}

\subsection{複数を作る}
集合$A$のすべての要素が集合$B$の要素でもあることを，\ko{$A$ skon $B$}と書く．「$A$は$B$の部分集合である」という意味で，$A=B$でも成り立つ．これを使って，「りんごだけからなる集合」を条件として作れる．
\begin{exe}
    \ex \ko{les leta skon kalamim fol}\hfill \sentence{りんごからなる集合}
\end{exe}
\ko{leta}は，\ko{le}を付ける主要部を省いた形である．この例では「何か」の位置を空け，りんごの集合の部分集合になるものを集めている．\ko{kalamle}という\emph{内複数}の形は，現行の暫定規約ではこの長形に対応する．項として使うときには，そこから場面に合う集合を取る．

これに対して\emph{外複数}の\ko{+na}は，ある個体を含む集合のクラスを作る．例えば\ko{cotona}は，場面から取った人を含む集合を表す．他の構成員まで全員人であることは要求しない．「私たち」のように使う場合には，そのクラスから文脈に合う集合を選ぶ．

内複数の長形には空集合や一個だけの集合も含まれる．複数という名前だけから「二個以上」とは読まない．個数を指定するには数詞を加える．例えば\ko{lu}は3，\ko{lulas kalamle}は三個のりんごからなる集合のクラスを表す．これだけでは，その集合が存在することも，その要素すべてが食べられることも述べていない．集合の条件，集合を選ぶこと，各要素について何かを述べることを，別々の操作として組み立てる．

\section{性質・量・比較}
「そのりんごは赤い」と「そのペンは長い」は，日本語では同じような形をしている．しかし，赤という色を持つことと，ある基準より長いことでは，述べている内容が違う．コノメノでは，この違いを性質の捉え方から整理する．

\subsection{値を持つこと}
\ko{fylam}は赤という色を，\ko{ansa}は値とその担い手の関係を表す．第1項が値，第2項がその値を持つ対象である．
\begin{exe}
    \ex \ko{fylam ansa kalam}\hfill \sentence{そのりんごは赤い}
    \ex \ko{kwilis ansa litof}\hfill \sentence{そのペンは黒い}
\end{exe}
\ko{kwilis}は「黒」，\ko{litof}は「ペン」である．\ko{fylam}や\ko{kwilis}は，ここでは色の値を指す項として使われている．「赤い」は，一項述語をりんごの後ろに置く形ではなく，赤という値とりんごを\ko{ansa}で結ぶ形になる．

色だけでなく，長さや重さも値を持つ．例えばペンの全長について話しているなら，次の形を使える．本書では，語形をまだ導入していない内容を\kox{日本語}で代用する．角括弧の中は発音するコノメノの語ではない．
\begin{exe}
    \ex \kox{18cm} \ko{ansa litof}\hfill \sentence{そのペンの全長は18cmだ}
\end{exe}
色や単位付きの数値のように，値を指定して表す量を\emph{定値量}という．長さの値としては18cmと0.18mは同じであり，書き方の違いは値の違いにはならない．測定して得た数値・単位・誤差の記録について話したい場合には，その記録を別の対象として扱う．

\ko{ansa}だけでは，どの部位のどの方向を測るかまでは言っていない．同じ箱にも高さ・幅・奥行きがある．「長さ」という種類が決まっても，どの長さかはまだ決まらないので，上の例では全長の話であることを先に揃えていた．数値を足せば曖昧さがすべて消えるわけではない．

\subsection{基準に照らして評価すること}
18cmという値が分かっても，それが長いかどうかは別の問題である．ペンとしては長くても，杖としては短いかもしれない．「長い」「重い」のような評価には，値に加えて，対象や使い道に応じた基準が関わる．こうした量を\emph{形容量}という．

\ko{stoona}「長い」，\ko{kadoa}「重い」は，\ko{locka}で対象と結ぶ．
\begin{exe}
    \ex \ko{stoona locka litof}\hfill \sentence{そのペンは長い}
    \ex \ko{kadoa locka folka}\hfill \sentence{そのカゴは重い}
\end{exe}
\ko{locka}の第1項が評価する程度，第2項が評価される対象である．「長い」の場合は，その対象の長さが，長いとされる側に入ることを述べる．「長さがある」と「長い」は区別する必要がある．短いペンにも長さはあるからである．

この区別を日本語の形容詞の語尾から決めることはできない．コノメノでは，語ごとに何を比べ，どんな基準を使うかを調べて分類する．長さのような一つの尺度による評価もあれば，「良い」のように複数の観点が関わる評価もある．

\begin{remark}
    「ペンとして」「この用途には」のように基準を明示して渡す口語構文は，まだ整備の途中である．上の例では基準を文脈で共有している．基準に依存するという意味上の説明と，任意の基準を短く言えることは区別する．
\end{remark}

程度をさらに評価することもできる．\ko{mofloa}は「とても」に当たり，\ko{locto}で程度と結ぶ．
\begin{exe}
    \ex \ko{mofloa locto stoona}\hfill \sentence{とても長い}
\end{exe}
ここで評価されるのは長さの程度である．りんごが赤という値を持つこと，ペンの長さを長いと評価すること，その長さの程度を「とても」と評価することを，それぞれ\ko{ansa}・\ko{locka}・\ko{locto}で書き分ける．

\subsection{二つの値を比較する}
「長い」は基準との比較だったが，「こちらの方が長い」は二つの値の比較である．短いペン同士にも，どちらが長いかという違いはある．逆に，長いと評された二本が同じ長さとは限らない．

値の順序を表す\ko{liimes}は「第1項は第2項以下である」，\ko{sammes}は「第1項は第2項以上である」を意味する．二本のペンの全長をそれぞれ$p,q$と書くことにすれば，次のように表せる．$p,q$は説明用の記号であり，実際には比較する値を指す項を置く．
\begin{exe}
    \ex \ko{$p$ liimes $q$}\hfill \sentence{$p$は$q$以下だ}
    \ex \ko{$p$ sammes $q$}\hfill \sentence{$p$は$q$以上だ}
\end{exe}
「以下」「以上」は等しい場合も含む．例えば$p=q=18$cmなら，この二文は両方とも成り立つ．したがって，\ko{liimes}だけを「より短い」と訳してはいけない．厳密に短いと言いたいなら，$p$が$q$以下で，しかも同じ値ではないことを述べる．
\begin{exe}
    \ex \ko{$p$ liimes $q$, meide $p$ kast $q$ tos}
    \glt \sentence{$p$は$q$より小さい値だ}
\end{exe}
この例の$p,q$は各出現で同じ値を指す．「長さ」を表す裸の概念語を二度言うだけでは，その同一性を指定したことにはならない．

比較するのは同じ種類の値でなければならない．18cmと20cmなら長さの順序を使えるが，18cmと20kgをそのまま比べる順序はない．「私は小町より背が低い」も，まず二人から身長の値を取り出して，その値を比べるという構成になる．対象を直接並べるだけの短い比較形は未整備なので，本書では値を取り出す構成までを押さえる．

\subsection{性質を文の中で使う}
性質表現も，前に学んだ関係節や量化と組み合わせられる．例えば「赤いりんご」は，次のように条件を作って表す．
\begin{exe}
    \ex \ko{les fylam ansa le kalam fol}\hfill \sentence{赤いりんご}
    \ex \ko{fylam ansa kalam kon}\hfill \sentence{赤いりんごがある}
\end{exe}
最初は赤いりんごの条件を作り，後はその条件を満たすりんごが存在すると述べる．どちらも\ko{ansa}の意味は同じである．性質の種類を語彙で選び，そこから文を組み立てる部分は既習の規則を使う．この分担によって，文法の共通性を保ちながら，個々の性質の違いを詳しく作り込める．

\section{引き下げと概念の分類}
食べる出来事を取り出したことで，食べる人と食べられるものを別々の役割で結べた．性質の文でも，同じように間に対象を置いてみると，短い関係語が何をまとめていたかが見えてくる．

\subsection{そのものの個別の質}
二つのりんごが同じ赤色でも，一方のりんごの赤さと，他方のりんごの赤さは区別できる．一つが変色しても，もう一つまで変色したことにはならない．このような，特定の対象に内在する個別の質を\emph{トロープ}といい，\ko{canon}で表す．色の値そのものと，その値を持つ個別の質は別の対象である．

トロープには，値と担い手が結び付く．\ko{nila}は値を第1項，トロープを第2項に取り，\ko{noccea}は担い手を第1項，トロープを第2項に取る．りんごが赤いという場合には，次の図になる．
\begin{center}
\begin{tikzpicture}[>=Stealth]
    \node (v) at (0,0) {赤という値};
    \node (c) at (3.2,0) {個別の赤さ};
    \node (x) at (6.4,0) {りんご};
    \draw[->] (v) -- node[above]{\ko{nila}} (c);
    \draw[->] (x) -- node[above]{\ko{noccea}} (c);
\end{tikzpicture}
\end{center}
左から一筆で読むなら，後半の\ko{noccea}を転置する．また，ここでは条件を満たす個別の質があると述べたいので，\ko{canon}には存在を表す\ko{kon}を付ける．
\begin{exe}
    \ex\begin{xlist}
        \ex \ko{fylam ansa kalam}\hfill 短い関係で述べる
        \ex \ko{fylam nila canon kon menoccea kalam}\hfill 個別の質を取り出す
    \end{xlist}
    \glt \sentence{そのりんごは赤い}
\end{exe}
後の形は，「赤という値を持ち，そのりんごに内在する個別の質がある」と読む．出来事を中央に置いた食事の図と同じく，二つの関係が中央の対象を共有している．短い\ko{ansa}は，こうした質を介する関係をまとめていたのである．

\begin{remark}
    コノメノの現行の捉え方では，値が変われば別のトロープになる．りんごが赤から別の色に変わることを，同じトロープが別の値を持つこととしては扱わない．これは性質のモデルの選択であり，日常語の「同じ赤さ」だけから決まることではない．
\end{remark}

\subsection{引き下げの共通形}
二つの対象$x,y$の関係$R$を，間に置いた対象$e$と，そこへ向かう二つの関係$K,L$で表す対応を\emph{引き下げ対応}という．条件を一つの式にまとめると，次の形になる．
\[
    xRy\quad\Longleftrightarrow\quad
    \exists e\in E\;\bigl(xKe\land yLe\bigr).
\]
$E$は間に置く対象の範囲，$\exists$は「存在する」，$\land$は「かつ」，$\Longleftrightarrow$は「両側が同じ条件を表す」という意味である．つまり「$x$と$y$が$R$で結ばれるのは，両方を役割$K,L$で結べる$e$があるとき，かつそのときに限る」と読む．

りんごの例では，$R$が\ko{ansa}，$E$がトロープの範囲，$K$が\ko{nila}，$L$が\ko{noccea}に当たる．この式では，何か適当な対象を挿入しただけでは足りない．間の対象を取り出した形から元の関係が得られ，元の関係があればそのような対象を取れる，という両方向の対応が必要である．

間の対象を明示すると，それについてさらに述べられる．例えば「長い」を精しく書くなら，ペンの長さを担うトロープを取り出し，その値が長いと評価される範囲に入ることを\ko{lofna}で表す．
\begin{exe}
    \ex\begin{xlist}
        \ex \ko{stoona locka litof}
        \ex \ko{stoona lofna canon kon menoccea litof}
    \end{xlist}
    \glt \sentence{そのペンは長い}
\end{exe}
赤さの例との違いは，値を指定する\ko{nila}が，評価する範囲との関係\ko{lofna}に替わった点にある．いずれも担い手への関係は\ko{noccea}である．\ko{ansa}と\ko{locka}の違いを残したまま，共通の仕組みで詳しく書ける．

\subsection{何を対象として立てるか}
引き下げの式だけでは，間の対象が出来事なのか，個別の質なのかは決まらない．それぞれについて，何が同じ対象で，どの関係が成り立ち，何がそこから言えるかを定める必要がある．このような概念と関係の体系を\emph{オントロジー}という．

本書に出てきた対象だけでも，りんごや人のような\emph{持続物}，食事のような\emph{生起物}，そのりんごの赤さのような\emph{個別従属物}，色や長さの値のような\emph{量}がある．物と出来事を区別していたから，意味役割の向きや省略された項も判断できた．概念の分類は，辞書を見やすくするだけでなく，文の読み方にも使われている．

ただし，「種類が同じ」と「同じ役割を持つ」も別である．りんごは，食事では食べられるものになり，所有関係では所有されるものになる．食べられる役割を持ったからといって，りんごという種類が役割に置き換わるわけではない．

コノメノでは，共通する原理を見つけることと，異なるものの立ち位置を分類することの両方を制作上の関心としている．引き下げは表現を共通の形にまとめ，オントロジーはその中で使う対象の違いを作り込む．すべてを同じ対象として扱えば簡単になるが，その区別が必要な文を表せなくなる．反対に，表現ごとに独立した仕組みを増やせば，語を替えて使うことが難しくなる．

分類やモデル化は一意には決まらない．現行の定義に沿って使う一方で，表したい例がうまく収まらなければ，その例を手掛かりに理論を改める．『詳説』が大きくなるのは，こうした個々の表現についての検討も蓄積しているためである．

\section{複文と発話}
文を続けるときには，「両方が成り立つ」と言うだけの場合もあれば，二つの出来事の時間や因果関係を述べる場合もある．さらに，文の内容そのものについて話したり，答えを相手に求めたりすることもある．これらを区別して組み立てる．

\subsection{文の論理結合と出来事の関係}
二つの文がともに成り立つことを表すのが\ko{lite}「かつ」である．これはしばしば省かれ，文を並べて表す．どちらか一方または両方が成り立つことは\ko{liko}「または」で表す．
\begin{exe}
    \ex \ko{komatci coto lite waka coto}\hfill \sentence{小町は人であり，私も人である}
    \ex \ko{komatci coto liko waka coto}
    \glt \sentence{小町が人であるか，私が人であるか，または両方だ}
\end{exe}
\ko{liko}は「両方ではない」とは言わない．また，\ko{lite}で二つの出来事を述べても，どちらが先か，一方が他方の原因かは，その結合だけからは決まらない．

出来事の間の関係は，出来事を項として結ぶ．例えば\ko{na}は同時の関係，\ko{noolja}は移動を表すので，次のように言える．
\begin{exe}
    \ex \ko{hamin na noolja}\hfill \sentence{食事と移動が同時に起こる}
\end{exe}
この文は，食べる出来事と移動する出来事を\ko{na}で結んでいる．両方の参与者が同じかどうかは，さらにその人をそれぞれの出来事に結んで示す．二つを一文にしただけで，同じ人の食事と移動になるわけではない．

順接の\ko{suinot}，逆接の\ko{filnot}，原因と結果を結ぶ\ko{luis}なども，状況を結ぶ関係である．これらを使うときにも，関係の両端を確認する．食事の文の末尾がりんごなら，そのまま次の関係語を足すとりんごから続く．食事を接続先にしたければ，指示詞などで食事へ戻る必要がある．

\subsection{文の内容を対象にする}
「そのりんごは私のものだ」と述べることと，「『そのりんごは私のものだ』という文を覚えている」と述べることは違う．後の場合には，文を使ってりんごの所有者を述べる代わりに，文そのものを話題にしている．

表現を引用する括弧は\ko{min}と\ko{tos}である．
\begin{exe}
    \ex \ko{min kalam micto waka tos}
    \glt \sentence{「そのりんごは私のものだ」という表現}
\end{exe}
これは中の文が真だと主張する形ではない．同じことを違う構文で言っても，引用した表現が同じとは限らない．文字列そのものを引用する\ko{minac}と\ko{tos}もあり，表現の構造を扱うことと，文字の並びを扱うことを区別できる．

\begin{remark}
    引用の形式的な定義は，引用対象を一つだけ含むクラスを返す．引用個体を項として使うときには，その一つの対象を取り出す．引用の括弧と個体の取り出しを，同じ操作としては扱わない．本書では，パラメータを含む引用の一般則までは扱わない．
\end{remark}

表現の形より粗く，どのような場合に文が成り立つかをまとめて扱うこともできる．コノメノでは，文が成り立つ世界のクラスを\emph{命題}として扱う．ここでいう世界は，語の解釈を与える形式的な対象であり，話し合っている可能な場合を表す．

命題を作る括弧は\ko{fol}と\ko{sus}である．\ko{sus}は命題化の冠詞\ko{sun}の括弧端形に当たる．りんごと話者を固定して考えると，次の句は，そのりんごを話者が所有する世界のクラスを作る．
\begin{exe}
    \ex \ko{fol kalam micto waka sus}
    \glt \sentence{そのりんごが私のものであること}
\end{exe}
同じ真偽条件を持つ文は，同じ世界のクラスを作ることがある．しかし，そのことから二つの表現も同じとは言えない．例えば，計算式の形について話したいのか，計算結果について話したいのかでは，残すべき区別が違う．何を同じとみなしたいかに応じて，表現の引用と命題化を使い分ける．

\subsection{成立の仕方を指す}
文を成り立たせる具体的な状況について話す場合もある．二つのりんごがそれぞれ赤いとき，二つの文はともに真でも，それを支える個別の赤さは違う．引き下げで取り出した対象は，この成立の仕方を表す手掛かりになる．

『詳説』では，文を正確に成り立たせる状況のクラスを作る冠詞を$C$と書く．$C$は本節での説明用の記号であり，そのまま発音する語ではない．クラスのうちどの状況を指すかは，さらに選ぶ必要がある．文の表現，その成立の仕方，成立する世界の範囲は，同じものではない．

この違いは，否定文をつなぐときに表面化する．「食べなかったので移動した」の前半を，食べる出来事の項で代用すると，起こらなかったはずの食事を置いてしまう．否定内容を支える状況を取る必要があるが，負の状況をどう与えるかと，その口語の長形には未整備の部分がある．肯定の出来事を結ぶ例から，否定前件の構文をそのまま作れるとはしない．

また，引用や命題化で作った内容を「知る」「望む」などに結ぶには，その動詞がどの種類の内容を取るかを確認する．日本語で「こと」と言えるものを全部同じ形にしない．複数の意味モデルを組み込むという制作意図は，こうした選び分けにも現れている．

\subsection{問いと答え}
平叙文をYes/No疑問にするには，\ko{hato}で始め，\ko{no}で答えを求める．
\begin{exe}
    \ex \ko{hato kalam micto non no}\hfill \sentence{そのりんごはあなたのものですか}
\end{exe}
答えは\ko{nai}「はい」，\ko{mei}「いいえ」で返せる．\ko{nai}は問いの中の内容を肯定し，\ko{mei}は否定する．これらは返答の慣用表現であって，「真」「偽」を表す論理定数とは役割が違う．

誰か，何かを尋ねる場合には，概念の前に疑問冠詞\ko{ena}を置く．\ko{ena coto}は「誰」，\ko{ena alkono}は「何」に当たり，それぞれ\ko{eto}，\ko{enata}という短い形もある．\ko{alkono}はここでは対象全般を範囲とする．
\begin{exe}
    \ex \ko{ha kalam micto ena coto no}\hfill \sentence{そのりんごは誰のものですか}
\end{exe}
この文は\ko{ha}で始め，所有者の位置を尋ねている．\ko{ena coto}に代えて\ko{eto}を使ってもよい．答えるときは，尋ねられた位置に入る人を伝える．自分のものなら，\ko{waka}「私」が答えになる．

\ko{ha}と\ko{tos}は，問いの選択肢を作る範囲を囲む．上のように末尾の\ko{no}で範囲の終わりが分かる場合は，\ko{tos}を省ける．\ko{ha}が選択肢を作ることと，\ko{no}が相手に回答を求めることを分けて考えると，文の内容と発話の働きを区別できる．疑問の形式意味論との対応には暫定の部分があるが，本節の形は現行の実践編に従う．

命令や依頼については，通常の必須文法として教えられるまでの整備がまだ揃っていない．問いが状況によって依頼のように使われることと，命令の規則が定義されていることを区別する．

\section{短い文章を読む}
出かける前に，カゴの中や手元の持ち物を確認している場面を考えよう．下の文章には，これまでに出た語と構文だけを使っている．まず，指示詞がどの対象を受けるかを追いながら読んでみてほしい．

\begin{quote}
\ko{kalam litam folka.}\\
\ko{tata micto waka.}\\
\ko{tan kalam ansa fylam.}\\
\ko{litof micto waka.}\\
\ko{tan litof ansa kwilis.}\\
\ko{tan litof locka stoona.}
\end{quote}

\paragraph{訳} そのりんごはそのカゴに入っている．そのりんごは私のものだ．そのりんごは赤い．そのペンは私のものだ．そのペンは黒い．そのペンは長い．

\subsection{対象を追う}
一文目では，りんご，カゴの順に項が現れる．文末ではカゴが最も新しいので，二文目の\ko{tata}は二つ前のりんごを受ける．りんごへ戻ってから私に言及するため，二文目の終わりでは，私，りんご，カゴの順になる．したがって，三文目では\ko{tan kalam}でりんごを受けられる．

四文目では新しくペンを導入し，その後で私に言及する．このとき私を新しい人として追加するのではなく，同じ話者への言及を先頭へ移す．五文目の\ko{tan litof}は，その次に新しいペンを受ける．五文目の終わりでは黒という値が先頭に来るので，六文目でもペンは\ko{tan litof}になる．色や程度の値も項として言及されれば，指示詞が数える対象に入る．

三・五文目の\ko{ansa}と六文目の\ko{locka}は，元の向きと逆に読んでいる．前者は「色の値・関係・担い手」，後者は「程度・関係・担い手」が元の順だが，この文章では指示詞で担い手に戻り，そこから色や程度へ進んでいる．対象の種類が合う向きに関係を読むことで，同じ構成を違う順に使える．

\subsection{少し書き替える}
次の二つを試してみよう．それぞれ，後続の文を伴わない独立した書き替えとして考える．
\begin{enumerate}
    \item りんごについて「私のものではない」と言う．
    \item 新たにペンを指し，「私のものだ」という内容をYes/No疑問にする．
\end{enumerate}
一つ目は\ko{meide kalam micto waka tos}である．すでに選んだりんごを再指示するなら，その時点の言及順に合わせて\ko{kalam}を指示詞に替える．二つ目は\ko{hato litof micto waka no}となる．相手の所有物か尋ねたい場合は，所有者も\ko{non}に替える．

文を直すときは，語だけでなく後続の照応も確認する．例えば三文目を色の値から始めれば，その後の言及順が変わる．一つ一つの文が文法的でも，文章全体では指示先がずれることがある．

\section{『詳説』への案内}
短い文を組み立てるときは，対象を決め，関係を選び，必要なら冠詞や括弧を付ける．その先で迷う場所は，用途によって違う．『詳説』では，次の箇所が入口になる．
\begin{itemize}
    \item 発音や接辞の形を確認する：音韻章「現行の音韻・語形体系」．独自文字は文字章「マリータスユ」．
    \item 冠詞・関係節・量化を確認する：実践編「表現の構成」．形式的な定義は形式文法章「基本体系」．
    \item 語順・省略・指示詞を確認する：実践編「文型と照応」，形式文法章「線形文モジュール」「指示詞モジュール」．
    \item 動詞にどの意味役割を使うか調べる：実践編「各種の表現」と，生起物の章の各概念・役割の定義．
    \item 性質や比較の捉え方を調べる：量・性質モジュール．所有・部分・内部の違いは持続物の章．
    \item 引き下げ・引用・命題化を調べる：形式意味論章．発話と照応の更新は形式文法章「談話モジュール」．
\end{itemize}

辞書では，日本語訳だけでなく，その語の上位概念と引数を見る．例えば\ko{micto}なら，「所有」という訳に加えて，第1項が所有されるもの，第2項が所有者であることを確かめる．\ko{hamin}のような概念語では，どんな出来事の一種か，参与者にどの役割を使うかも調べる．語彙・冠詞・音列・慣用表現・括弧は辞書内で区分されており，同じ綴りでも同じ働きとは限らない．

本書をLLM共同執筆の足場にする場合も，実際の修正前には『詳説』の対象箇所と依存箇所へ戻る．本書にない規則を「存在しない」とみなしたり，省略した説明を自由な解釈の余地として埋めたりしない．例を作って収まりが悪ければ，語彙の不足か，構文の未整備か，概念の分類やモデルの問題かを切り分ける．

理論上の定義と検討中の案は『詳説』本文が基準になる．辞書データは実装済み語彙の基準であり，本文の改訂がまだ移っていない箇所もある．本書にも，暫定の規則はその近くに記した．新しい例や説明から理論を改める場合には，『素描』だけに例外を増やさず，『詳説』の定義と依存する説明を一緒に直す．
\end{document}
